
\documentclass[11pt,a4paper]{article}

\usepackage[utf8]{inputenc}
\usepackage[T1]{fontenc}
\usepackage[brazil]{babel}
\usepackage{lmodern}
\usepackage{geometry}
\usepackage{amsmath,amssymb}
\usepackage{booktabs}
\usepackage{enumitem}
\usepackage{listings}
\usepackage{hyperref}
\usepackage{microtype}
\usepackage{array}
\usepackage{fancyhdr}

\geometry{margin=2.5cm}

\hypersetup{
    colorlinks=true,
    linkcolor=black,
    urlcolor=blue
}

\pagestyle{fancy}
\fancyhf{}
\lhead{Análise e Projeto de Algoritmos}
\rhead{Trabalho Prático}
\cfoot{\thepage}

\setlist[itemize]{topsep=3pt,itemsep=2pt,parsep=0pt}
\setlist[enumerate]{topsep=3pt,itemsep=2pt,parsep=0pt}

\lstset{
    basicstyle=\ttfamily\small,
    frame=single,
    columns=fullflexible,
    keepspaces=true,
    showstringspaces=false,
    breaklines=true
}

\begin{document}

\begin{center}
    {\Large \textbf{Análise e Projeto de Algoritmos}}\\[0.5em]
    {\large \textbf{Trabalho Prático}}\\[0.5em]
    {\LARGE \textbf{Seleção de Postos Centrais de Monitoramento}}
\end{center}

\section{Contexto}

Um parque florestal possui postos de monitoramento distribuídos em sua área. Cada posto conta com câmeras, sensores e equipamentos de comunicação.

Devido à vegetação, ao relevo e à distância entre os equipamentos, nem todos os postos conseguem comunicar-se diretamente. Quando dois postos possuem comunicação direta, eles podem compartilhar alertas e informações.

Alguns postos serão escolhidos como \textbf{centrais de monitoramento}. Uma central é responsável por receber, consolidar e encaminhar os alertas dos postos diretamente conectados a ela.

Cada posto deve satisfazer pelo menos uma das seguintes condições:

\begin{itemize}
    \item ser escolhido como central; ou
    \item possuir comunicação direta com pelo menos uma central.
\end{itemize}

O objetivo é escolher a menor quantidade possível de postos centrais.

\section{Representação}

A rede será representada por um grafo simples e não direcionado
\[
G=(V,E),
\]
em que cada vértice representa um posto e cada aresta representa uma comunicação direta.

Uma solução é um conjunto
\[
D\subseteq V
\]
tal que, para todo vértice $v\in V$,
\[
v\in D
\]
ou existe $u\in D$ tal que
\[
\{u,v\}\in E.
\]

O objetivo é minimizar
\[
|D|.
\]

\section{Parte A -- Parque geral}

Nesta parte, a disposição dos postos é arbitrária.

O grupo deverá desenvolver e implementar um \textbf{algoritmo exato} para encontrar uma solução de cardinalidade mínima.

Um algoritmo exato deve:

\begin{enumerate}
    \item produzir uma solução válida;
    \item garantir que nenhuma solução menor existe;
    \item terminar para toda entrada válida, ainda que o tempo de execução possa ser elevado.
\end{enumerate}

Não será indicada uma técnica específica. O grupo deverá decidir como representar e examinar as possíveis soluções.

\subsection{Entrada}

O programa deverá aceitar o formato DIMACS:

\begin{lstlisting}
c comentario
p edge n m
e u1 v1
e u2 v2
...
e um vm
\end{lstlisting}

Os vértices serão identificados por
\[
1,2,\ldots,n.
\]

O grafo poderá ser desconexo.

\subsection{Saída}

O programa deverá apresentar:

\begin{itemize}
    \item o tamanho da melhor solução encontrada;
    \item os postos escolhidos;
    \item se a otimalidade foi comprovada;
    \item o tempo de execução;
    \item o motivo da interrupção, quando aplicável.
\end{itemize}

Exemplo:

\begin{lstlisting}
Quantidade minima de centrais: 2
Postos escolhidos: 1 4
Otimalidade comprovada: sim
Tempo de execucao: 0.00231 segundos
\end{lstlisting}

Caso o limite de tempo seja atingido:

\begin{lstlisting}
Melhor solucao encontrada: 37
Otimalidade comprovada: nao
Tempo de execucao: 600.000 segundos
Motivo da interrupcao: limite de tempo
\end{lstlisting}

\subsection{Testes da Parte A}

Todos os grupos deverão utilizar o mesmo conjunto de instâncias fornecido pela professora.

As instâncias incluem:

\begin{itemize}
    \item grafos completos;
    \item grafos sem arestas;
    \item estrelas;
    \item caminhos;
    \item ciclos;
    \item grafos desconexos;
    \item grafos aleatórios;
    \item grafos esparsos e aproximadamente regulares;
    \item uma instância grande para observar os limites práticos da solução exata.
\end{itemize}

Não se espera que todas as instâncias grandes sejam resolvidas com otimalidade comprovada.

\section{Parte B -- Parque linear}

Considere agora um segundo parque, longo e estreito, com uma única trilha principal.

Cada posto monitora um trecho contínuo da trilha, determinado por duas posições:
\[
\ell_i \leq r_i.
\]

Dois postos comunicam-se diretamente quando existe pelo menos uma posição monitorada por ambos.

O objetivo continua sendo escolher o menor número possível de centrais.

\subsection{Atividades da Parte B}

O grupo deverá:

\begin{enumerate}
    \item aplicar o algoritmo da Parte A às instâncias do parque linear;
    \item analisar quais propriedades adicionais aparecem devido à disposição dos postos;
    \item investigar um algoritmo exato que explore diretamente essa estrutura;
    \item comparar o novo algoritmo com o algoritmo geral.
\end{enumerate}

Como orientação, o grupo poderá investigar:

\begin{itemize}
    \item projeto por indução;
    \item programação dinâmica;
    \item estratégia gulosa.
\end{itemize}

Não é necessário implementar uma solução baseada em cada paradigma. O grupo deverá escolher e justificar a abordagem considerada mais adequada.

\subsection{Requisitos teóricos}

Para o algoritmo desenvolvido na Parte B, o grupo deverá apresentar:

\begin{itemize}
    \item pseudocódigo;
    \item prova de que todos os postos são atendidos;
    \item prova de que a solução é mínima;
    \item análise de complexidade;
    \item pelo menos um contraexemplo para uma estratégia intuitiva que não produza sempre uma solução ótima.
\end{itemize}

A análise deverá responder se o novo algoritmo permanece exponencial ou se passa a ser polinomial.

\subsection{Testes da Parte B}

Cada instância será fornecida em dois formatos:

\begin{itemize}
    \item arquivo \texttt{.col}, contendo o grafo;
    \item arquivo \texttt{.intervalos.txt}, contendo os extremos dos trechos.
\end{itemize}

O formato do segundo arquivo será:

\begin{lstlisting}
posto inicio fim
\end{lstlisting}

\section{Protocolo de medição}

O limite máximo de cada execução será de
\[
600\text{ segundos}.
\]

O algoritmo deverá ser executado de forma sequencial, utilizando apenas uma thread.

Antes das medições, o grupo deverá:

\begin{itemize}
    \item encerrar programas desnecessários;
    \item não executar diferentes instâncias simultaneamente;
    \item não realizar outras tarefas durante a execução;
    \item utilizar a mesma máquina e as mesmas opções de compilação nos testes comparados;
    \item manter o computador conectado à energia elétrica;
    \item evitar modos de economia de energia.
\end{itemize}

Não será exigida a limpeza manual dos caches do processador ou do sistema operacional.

Cada instância concluída deverá possuir uma execução de aquecimento, não registrada, seguida de pelo menos três execuções oficiais.

O tempo deverá ser medido com relógio monotônico de alta resolução, por exemplo:

\begin{itemize}
    \item \texttt{std::chrono::steady\_clock}, em C++;
    \item \texttt{System.nanoTime()}, em Java;
    \item \texttt{time.perf\_counter()}, em Python.
\end{itemize}

O valor principal apresentado deverá ser a mediana das execuções.

Instâncias que atinjam 600 segundos não precisam ser repetidas.

\section{Relatório}

O relatório deverá conter:

\begin{enumerate}
    \item modelagem do problema;
    \item descrição e pseudocódigo do algoritmo da Parte A;
    \item prova de correção da Parte A;
    \item análise de tempo e espaço da Parte A;
    \item resultados experimentais da Parte A;
    \item descrição da estrutura do parque linear;
    \item algoritmo proposto para a Parte B;
    \item prova de correção e de otimalidade da Parte B;
    \item análise de complexidade da Parte B;
    \item comparação experimental entre os algoritmos;
    \item dificuldades, erros e decisões tomadas durante o desenvolvimento.
\end{enumerate}

O relatório poderá ter entre 6 e 10 páginas, sem contar capa, referências e anexos.

\section{Arquivos a entregar}

A entrega deverá conter:

\begin{enumerate}
    \item código-fonte;
    \item instruções de compilação e execução;
    \item relatório em PDF;
    \item arquivo com os resultados brutos;
    \item identificação dos integrantes.
\end{enumerate}

Não será permitido utilizar bibliotecas que resolvam diretamente o problema de otimização.

\section{Critérios de avaliação}

\begin{center}
\begin{tabular}{>{\raggedright\arraybackslash}p{10cm}r}
\toprule
\textbf{Critério} & \textbf{Peso} \\
\midrule
Correção e exatidão do algoritmo geral & 25\% \\
Prova e análise de complexidade da Parte A & 15\% \\
Análise experimental da Parte A & 15\% \\
Algoritmo especializado para o parque linear & 20\% \\
Prova e análise de complexidade da Parte B & 15\% \\
Organização, clareza e reprodutibilidade & 10\% \\
\bottomrule
\end{tabular}
\end{center}

\section{Observações finais}

\begin{itemize}
    \item Toda fonte consultada deverá ser citada.
    \item O grupo deverá compreender integralmente o código entregue.
    \item Qualquer integrante poderá ser solicitado a explicar uma parte do algoritmo ou justificar a resposta produzida.
    \item Encontrar uma solução válida não é suficiente: quando uma solução for declarada ótima, sua otimalidade deverá estar comprovada.
\end{itemize}

\end{document}
